The book of Section \ref{sec.theBook} is a snapshot.
This section describes how that snapshot changes.

Exactly three kinds of event modify a limit order book:
the \emph{submission} of a new limit order,
the \emph{cancellation} of a resting one,
and the \emph{execution} of a resting one against an incoming market order.
Everything a venue publishes is a sequence of such events,
and a model of the book is a model of their arrival.

It is natural to describe these arrivals with counting processes.
Fix a price level and a side, say the bid.
Let $\bidArrivals$ count the units of volume added at that level up to time $t$,
by submission,
and let $\bidDepartures$ count the units removed,
by cancellation or execution.
The queue at that level is then
\begin{equation}
\label{eq.queueDynamics}
\bidQueue = \bidQueue[0] + \bidArrivals - \bidDepartures,
\end{equation}
and the book is the collection of such queues across price levels.

Equation \eqref{eq.queueDynamics} is trivial as an accounting identity,
and it becomes a model as soon as one specifies the law of $\arrivals$ and $\departures$.
The simplest choice takes them to be independent Poisson processes with level-dependent
rates,
as in \cite{CST10sto}.
The resulting model is analytically tractable and reproduces the shape of the average book,
but it misses a feature that is impossible to ignore in real data:
order flow is strongly \emph{clustered} in time,
and events of one type make events of other types more likely.
Capturing that requires self- and cross-exciting processes,
whose intensity $\intensity$ depends on the history of the flow itself.

\begin{remark}
\label{remark.eventsAreNotIndependent}
The empirical autocorrelation of the signs of market orders
decays slowly, over hours rather than seconds.
This is not evidence of a predictable price:
it is mostly the footprint of large trades being split into many small ones.
Section \ref{sec.lobConclusions} returns to this point.
\end{remark}

A convenient one-number summary of the flow is the \emph{order flow imbalance}.
Over a time interval $(s,t]$ it aggregates the signed changes at the best quotes,
\begin{equation}
\label{eq.orderFlowImbalance}
\OFI_{s,t} :=
\big(\bidArrivals[t] - \bidArrivals[s]\big)
- \big(\bidDepartures[t] - \bidDepartures[s]\big)
- \big(\askArrivals[t] - \askArrivals[s]\big)
+ \big(\askDepartures[t] - \askDepartures[s]\big),
\end{equation}
counting as positive whatever pushes the price up and as negative whatever pushes it down.
\cite{CKS14pri} show that contemporaneous mid-price changes are close to linear in
$\OFI_{s,t}$,
with a slope that scales with the inverse of the depth.
This is one of the cleanest empirical regularities in microstructure,
and it is a good first exercise on real data.

\subsection{Price impact}
\label{sec.priceImpact}

A trade moves the price against the trader who initiated it.
This is \emph{price impact},
and it is the dominant cost of trading at institutional size.

The mechanism is visible directly in the book.
A market order to buy consumes the ask queue from the best price upwards,
so the next unit is bought at a worse price than the previous one,
and after the order has passed the best ask has moved up.
Part of that move persists,
because other participants read the trade as information;
part of it decays,
as new limit orders replenish the consumed level.

Impact is the reason a large order cannot simply be sent to the market.
It is instead split into many small \emph{child} orders,
released over a period of time,
which is the subject of optimal execution;
\cite{AC01opt} is the classical reference,
and \cite{CJP15alg} gives the modern treatment.
The trade-off is exactly the one identified in
Section \ref{sec.orderDrivenMarkets}:
trading slowly reduces impact but increases exposure to price risk over the execution
window.
